(1) Gauss定律令 \(D_r(r)=Q/(2\pi Lr)\)。因为 \(dD/dE=\varepsilon+3\chi E^2>0\)，每个 \(D_r\) 恰有一个同号实解 \(E_r\)，不可能出现该简化材料的多稳态静电支。

(2) 电势差 \(V(Q)=\int_a^bE[Q/(2\pi Lr)]dr\)。当 \(\chi D^2/\varepsilon^3\ll1\)，反求 \(E=D/\varepsilon-\chi D^3/\varepsilon^4+O(\chi^2)\)，从而 \(V=\frac{Q\ln(b/a)}{2\pi L\varepsilon}-\frac{\chi Q^3}{2(2\pi L)^3\varepsilon^4}(a^{-2}-b^{-2})+O(\chi^2)\)；正的 \(\chi\) 降低固定电荷下的电压。

(3) 准静态可逆充电的电功密度是 \(E\,dD\)，故 \(u=\int_0^D E(D')dD'=\int_0^E E'(\varepsilon+3\chi E'^2)dE'=\varepsilon E^2/2+3\chi E^4/4\)。总能 \(U=2\pi L\int_a^br\,u(r)dr\)；由链式求导 \(dU/dQ=\int_a^bE_r(r)dr=V\)，符合电源功账本。

(4) 对材料方程求导 \(dE/dD=1/(\varepsilon+3\chi E^2)\)，因 \(D=Q/(2\pi Lr)\)，有 \(\boxed{C_{\rm diff}^{-1}=\frac1{2\pi L}\int_a^b\frac{dr}{r[\varepsilon+3\chi E_r(r)^2]}}\)，显然正且大于线性 \(C\)。强场时 \(E\sim(D/\chi)^{1/3}\)，所以 \(V\propto Q^{1/3}\)、\(C_{\rm diff}\propto Q^{2/3}\)，但真实介质最终会偏离三次本构。
